\chapter{Systemic Sclerosis (Scleroderma)}
\index{Systemic Sclerosis (Scleroderma)}
\label{sec:Systemic Sclerosis (Scleroderma)}

\section{Overview}
Systemic sclerosis is a chronic autoimmune connective tissue disease characterized by widespread vasculopathy, immune dysregulation, and progressive scarring of the skin and internal organs, with an incidence of approximately 1--2 per 100,000, a female-to-male ratio of 4--6:1, and peaking in the fifth decade. This autoimmune condition occurs when the immune system mistakenly attacks the body's own tissues. It follows a chronic course with periods of flare and remission. Autoimmune diseases result from an abnormal immune response where the body mistakenly attacks its own healthy tissues. These conditions are typically chronic, with alternating periods of flare and remission. They can affect a single organ or be systemic, involving multiple organ systems. Women are affected more frequently than men for most autoimmune conditions. The exact prevalence varies by condition, but collectively autoimmune diseases affect a significant portion of the population.
\section{Causes}
This condition happens because of excessive collagen deposition by activated fibroblasts, driven by endothelial injury, autoimmune inflammation, and transforming growth factor-beta (TGF-$\beta$) signaling. In autoimmune conditions, a combination of genetic susceptibility and environmental triggers initiates an inappropriate immune response against self-antigens. This leads to chronic inflammation and tissue damage in affected organs. A combination of genetic susceptibility and environmental triggers initiates the autoimmune process. Specific HLA genotypes are strongly associated with many autoimmune conditions. Environmental triggers may include infections, smoking, medications, and hormonal changes. Loss of self-tolerance leads to activation of autoreactive T cells and B cells producing autoantibodies. Molecular mimicry between pathogen antigens and self-antigens may trigger cross-reactive immune responses.
\section{Risk Factors}
Family history of autoimmune disease Female sex (higher prevalence) Specific HLA genotypes Epstein-Barr virus infection Cigarette smoking Certain medications (drug-induced autoimmunity) Hormonal factors (pregnancy, postpartum) Vitamin D deficiency Stress and psychological factors Geographic and ethnic background

\begin{itemize}[noitemsep]
  \item Family history of autoimmune disease
  \item Female sex (higher prevalence)
  \item Specific HLA genotypes
  \item Epstein-Barr virus infection
  \item Cigarette smoking
  \item Certain medications (drug-induced autoimmunity)
  \item Hormonal factors (pregnancy, postpartum)
  \item Vitamin D deficiency
\end{itemize}
\section{Signs and Symptoms}
\begin{itemize}[noitemsep]
  \item You may experience Raynaud's phenomenon (the most frequent initial manifestation, present in over 95\%), and three clinical variants exist: limited cutaneous SSc (lcSSc, skin thickening distal to elbows/knees, formerly CREST syndrome: Calcinosis cutis, Raynaud's phenomenon, Esophageal dysmotility, Sclerodactyly, Telangiectasias), diffuse cutaneous SSc (dcSSc, skin thickening proximal to elbows/knees and trunk, more rapid and severe internal organ involvement), and systemic sclerosis sine scleroderma (Raynaud's phenomenon and internal organ involvement without clinically apparent skin scarring). lung involvement is the leading cause of death: lung scarring (more common in dcSSc) presents with shortness of breath, cough, and restrictive lung disease
  \item Lung high blood pressure (more common in lcSSc) presents with exertional shortness of breath, right heart failure, and carries a poor outlook. Other features include gastroesophageal reflux, gastric antral vascular ectasia (watermelon stomach), telangiectasias, and myositis.
  \item Fatigue and general malaise
  \item Joint pain, swelling, and stiffness
  \item Skin rashes and lesions
  \item Low-grade fever
  \item Muscle weakness and pain
  \item Organ-specific symptoms based on affected system
  \item Weight changes (loss or gain)
  \item Dry eyes and dry mouth (Sjogren syndrome)
  \item Raynaud phenomenon (cold-induced color changes in fingers)
  \item Fluctuating symptoms with periods of flare and remission
\end{itemize}
\section{Progression and Stages}
Autoimmune diseases typically follow a relapsing-remitting course with unpredictable flares. Early disease may present with nonspecific symptoms before organ-specific manifestations develop. Chronic inflammation over time can lead to irreversible tissue damage and organ dysfunction. With appropriate treatment, many patients achieve remission or low disease activity states.
\section{Complications}
Scleroderma kidney crisis (SRC) is a life-threatening complication of dcSSc (incidence 5--10\%), characterized by cancerous high blood pressure, rapidly progressive kidney failure, microangiopathic hemolytic anemia, and thrombocytopenia.

\begin{itemize}[noitemsep]
  \item Irreversible organ damage from chronic inflammation
  \item Joint deformities and erosions in inflammatory arthritis
  \item Renal failure in lupus nephritis
  \item Pulmonary fibrosis or hypertension
  \item Cardiovascular complications from systemic inflammation
  \item Osteoporosis from chronic corticosteroid use
  \item Increased risk of infections from immunosuppressive therapy
  \item Lymphoma risk in some autoimmune conditions
\end{itemize}
\section{Diagnosis}
Doctors usually diagnose this based on your symptoms and a physical exam, supported by autoantibodies (anti-centromere in lcSSc, anti-Scl-70/anti-topoisomerase I in dcSSc, anti-RNA polymerase III associated with SRC risk), nailfold capillaroscopy, and organ-specific testing (PFTs, echocardiogram, high-resolution CT chest). Diagnosis is based on a combination of clinical features, laboratory findings (autoantibodies, inflammatory markers), and imaging studies. Classification criteria help standardize the diagnostic process. Clinical history and physical examination Autoantibody testing (ANA, RF, anti-CCP, anti-dsDNA, etc.) Inflammatory markers (ESR, CRP) Complete blood count and comprehensive metabolic panel Imaging studies (X-ray, MRI, ultrasound) of affected areas Biopsy of affected tissue for histological examination Classification criteria for specific autoimmune diseases Rheumatology consultation for suspected autoimmune disease Monitoring for organ involvement (renal, pulmonary, cardiac) Genetic testing for HLA associations when indicated

\begin{itemize}[noitemsep]
  \item Clinical history and physical examination
  \item Autoantibody testing (ANA, RF, anti-CCP, anti-dsDNA, etc.)
  \item Inflammatory markers (ESR, CRP)
  \item Complete blood count and comprehensive metabolic panel
  \item Imaging studies (X-ray, MRI, ultrasound) of affected areas
  \item Biopsy of affected tissue for histological examination
  \item Classification criteria for specific autoimmune diseases
  \item Rheumatology consultation for suspected autoimmune disease
\end{itemize}
\section{Prevention}
\begin{itemize}[noitemsep]
  \item No known prevention for most autoimmune diseases
  \item Avoid known triggers when possible
  \item Maintain a healthy lifestyle to support immune function
  \item Early recognition of symptoms for prompt diagnosis
  \item Regular medical check-ups for monitoring
  \item Adequate vitamin D levels through sun exposure or supplementation
\end{itemize}
\section{Treatment Options}
Treatment typically involves organ-specific and multidisciplinary: calcium channel blockers, prostacyclin analogs, and endothelin receptor antagonists for Raynaud's phenomenon and lung high blood pressure

\begin{itemize}[noitemsep]
  \item Proton pump inhibitors and prokinetics for GI involvement
  \item ACE inhibitors for scleroderma kidney crisis (first-line, which has transformed outcomes---SRC mortality has decreased from 80\% to 15--20\%)
  \item Mycophenolate, cyclophosphamide, and nintedanib for interstitial lung disease
  \item And methotrexate for skin and joint involvement.
  \item Nonsteroidal anti-inflammatory drugs (NSAIDs) for symptom relief
  \item Corticosteroids for rapid control of inflammation
  \item Disease-modifying antirheumatic drugs (DMARDs)
  \item Biologic therapies targeting specific immune pathways
  \item Immunosuppressive agents (azathioprine, mycophenolate, cyclophosphamide)
  \item Antimalarial drugs (hydroxychloroquine) for mild disease
  \item Physical and occupational therapy to maintain function
  \item Pain management for chronic pain
  \item Lifestyle modifications including diet and stress reduction
\end{itemize}
\section{Living with It}
\begin{itemize}[noitemsep]
  \item Take medications consistently as prescribed
  \item Identify and avoid personal flare triggers
  \item Maintain a balanced diet and regular gentle exercise
  \item Practice stress management techniques
  \item Join support groups for emotional support
  \item Work with a rheumatologist for ongoing disease management
  \item Monitor for medication side effects
  \item Plan activities around energy levels during flares
\end{itemize}
\section{Outlook}
The outlook varies from person to person, with the diffuse cutaneous variant having a worse outlook, and lung involvement being the leading cause of death. Autoimmune conditions are typically chronic and require long-term management. With appropriate treatment, many patients achieve good symptom control and maintain quality of life. Autoimmune diseases are generally chronic conditions requiring long-term management. Disease activity can fluctuate, requiring adjustments in treatment over time. Early diagnosis and treatment improve outcomes and prevent irreversible organ damage. Research continues to develop more targeted therapies with fewer side effects.
